\chapter*{Introduction}
\addcontentsline{toc}{chapter}{Introduction}

This solutions manual follows the same progression as the text itself. The early
chapters build the mathematical language needed for automata theory; the middle
chapters develop finite automata, regular languages, and equivalent formalisms;
and the later chapters move on to grammars, pushdown automata, Turing machines,
and decidability. The chapter-level themes are:

\begin{enumerate}
\item \textbf{Chapter 0: Preliminaries.} The book begins with the mathematical
tools used everywhere else: logic, sets, relations, functions, cardinality,
induction, recursion, and Backus-Naur Form.

\item \textbf{Chapter 1: Introduction and Basic Definitions.} This chapter
introduces alphabets, words, and finite automata, then develops examples,
implementation ideas, and practical applications.

\item \textbf{Chapter 2: Characterization of Finite Automaton Definable
Languages.} The central ideas are right congruences, Nerode's Theorem, and
pumping arguments that explain exactly which languages finite automata can
recognize.

\item \textbf{Chapter 3: Minimization of Finite Automata.} The emphasis here is
on homomorphisms, isomorphisms, and algorithms for reducing automata to their
essential states.

\item \textbf{Chapter 4: Nondeterministic Finite Automata.} This chapter shows
how nondeterminism extends the model, how NDFAs compare with DFAs, and how
lambda-transitions and circuit realizations fit into the picture.

\item \textbf{Chapter 5: Closure Properties.} The main theme is that regular
languages are stable under many operations, and the chapter develops the
constructions that prove those closure results.

\item \textbf{Chapter 6: Regular Expressions.} The book studies the algebra of
regular expressions, shows how regular expressions describe the same class of
languages as finite automata, and uses language equations to connect the two
views.

\item \textbf{Chapter 7: Finite-State Transducers.} Here automata become
input-output devices, with attention to their definitions, minimization, Moore
machines, and concrete applications such as circuit-style implementations.

\item \textbf{Chapter 8: Regular Grammars.} This chapter places regular
grammars in the larger grammar hierarchy and relates right-linear grammars to
automata and regular expressions.

\item \textbf{Chapter 9: Context-Free Grammars.} The focus shifts to parse
trees, ambiguity, canonical forms, pumping results, and closure properties for
context-free languages.

\item \textbf{Chapter 10: Pushdown Automata.} The main ideas are the PDA model,
its equivalence with context-free grammars, alternative acceptance conventions,
and the special behavior of deterministic PDAs.

\item \textbf{Chapter 11: Turing Machines.} This chapter develops the Turing
machine model, compares its variants, relates it to other language-defining
systems, and begins to measure computational power.

\item \textbf{Chapter 12: Decidability.} The final chapter studies what can and
cannot be decided algorithmically, starting with regular-language questions and
ending with undecidability and Turing-decidable languages.
\end{enumerate}
